\section{Билет 4. Закон движения для гармонических колебаний. Скорость и ускорение при гармоническом колебании. Графики смещения, скорости и ускорения.}

\begin{definition}
Закон движения гармонического осциллятора задаётся уравнением:
\begin{equation}
    x(t) = A \cos(\omega_0 t + \varphi).
\end{equation}
Кинематические характеристики (скорость и ускорение) находятся путём последовательного дифференцирования смещения по времени.
\end{definition}

\begin{theorem}[Скорость и ускорение]
Мгновенная скорость $v(t)$ равна первой производной смещения:
\begin{equation}
    v(t) = \dot{x}(t) = -A\omega_0 \sin(\omega_0 t + \varphi) = A\omega_0 \cos\left(\omega_0 t + \varphi + \frac{\pi}{2}\right). \label{eq:velocity}
\end{equation}
Ускорение $w(t)$ равна второй производной смещения:
\begin{equation}
    w(t) = \ddot{x}(t) = -A\omega_0^2 \cos(\omega_0 t + \varphi) = A\omega_0^2 \cos(\omega_0 t + \varphi + \pi). \label{eq:acceleration}
\end{equation}
\end{theorem}

\begin{property}
Из соотношений \eqref{eq:velocity} и \eqref{eq:acceleration} следуют ключевые свойства:
\begin{enumerate}
    \item Скорость и ускорение также изменяются по \textbf{гармоническому закону}.
    \item Ускорение и смещение находятся в \textbf{противофазе}: когда смещение достигает наибольшего положительного значения, ускорение достигает наибольшего отрицательного значения, и наоборот.
    \item Ускорение всегда имеет знак, противоположный знаку смещения, следовательно, сила, заставляющая тело совершать гармонические колебания, направлена всегда в сторону положения равновесия.
\end{enumerate}
\end{property}

\begin{remark}
\textbf{Графики смещения, скорости и ускорения.} При совместном построении графиков $x(t)$, $v(t)$ и $w(t)$ наблюдается следующее:
\begin{itemize}
    \item График скорости $v(t)$ опережает график смещения $x(t)$ по фазе на $\pi/2$.
    \item График ускорения $w(t)$ сдвинут относительно $x(t)$ на $\pi$ (полная противофаза).
    \item Все три кривые являются гармоническими, но имеют разные амплитуды: $A$ для смещения, $A\omega_0$ для скорости и $A\omega_0^2$ для ускорения.
\end{itemize}
Фазовые сдвиги наглядно демонстрируют, что максимальная скорость достигается при прохождении положения равновесия ($x=0$), а максимальное по модулю ускорение --- в точках максимального отклонения ($x=\pm A$).
\end{remark}

\begin{figure}[htbp]
    \centering
    \begin{tikzpicture}
    \begin{axis}[
        xlabel = {$t$},
        ylabel = {$x, v, w$},
        xmin = 0, xmax = 4,
        ymin = -4.5, ymax = 4.5,
        domain = 0:4,
        samples = 200,
        grid = both,
        grid style = {dashed, gray!20}, % Очень легкая сетка, чтобы не мешать
        width = 0.9\textwidth,
        height = 0.5\textheight,
        axis lines = center,
        legend pos = outer north east,
        legend style = {fill=none, draw=none}, % Прозрачная легенда
        xtick = {0, 1.57, 3.14},
        xticklabels = {$0$, $\frac{\pi}{2\omega_0}$, $\frac{\pi}{\omega_0}$},
        ytick = {-4, -2, 0, 2, 4},
        clip=false % Чтобы линии не обрезались на границах
    ]

    % 1. Смещение (синий) - Амплитуда 1
    \addplot [blue, thick, smooth] {cos(deg(2*x))};
    \addlegendentry{$x(t)$ (Смещение)}

    % 2. Скорость (красный) - Амплитуда 2 (опережает на pi/2)
    \addplot [red, thick, smooth] {-2*sin(deg(2*x))};
    \addlegendentry{$v(t)$ (Скорость)}

    % 3. Ускорение (зеленый) - Амплитуда 4 (противофаза с x)
    \addplot [green!60!black, thick, smooth] {-4*cos(deg(2*x))};
    \addlegendentry{$w(t)$ (Ускорение)}

    % Тонкие серые линии для акцента на фазах (не перекрывают графики)
    \draw[gray, dotted, thin] (axis cs:1.57, -4.5) -- (axis cs:1.57, 4.5);
    \draw[gray, dotted, thin] (axis cs:3.14, -4.5) -- (axis cs:3.14, 4.5);

    \end{axis}
    \end{tikzpicture}
    \caption{Графики гармонических колебаний. Скорость $v(t)$ опережает смещение $x(t)$ по фазе на $\pi/2$, ускорение $w(t)$ находится в противофазе со смещением ($\pi$).}
    \label{fig:harmonic_xyz_clean}
\end{figure}

